\section{结论}

多年累积的卫星海洋遥感大数据不仅能为自主水下航行器提供先验底图，更能生成一套经误差标定且独立于搜索过程的留出想象海洋，作为衡量机载计算价值的客观标尺。本文围绕东海赤潮应急巡测中走、深想与停下复测的资源分配问题，将48个随采样局地化集合同化更新的想象海洋划分为32个搜索用海洋与16个留出海洋，提出了出发前全任务预演方法及结合留出估值与两倍标准误判据的在线单步元推理方法。在3197天东海卫星数据构建的422个赤潮窗口、四种机载硬件及两种真值模型上的实验表明：留出评判在跨三个半数量级的推演预算下与真实航线价值偏差不超过0.005；开启高功耗芯片深想仅在基础算力过慢时有益，而停下复测始终得不偿失；出发前预演与在线留出估值无需历史任务即达到依赖历史任务方法的水平，并优于经典单步元推理以及除慢计算机主设置下贪心规则外的全部巡测规则。未来工作将把该方法拓展至三维斜压海流场与多航行器协同赤潮巡测研究。
